\documentclass[11pt]{article}
\usepackage{fontspec}
\usepackage[french]{babel}
\usepackage{amsmath}
\usepackage{amssymb}
\usepackage{amsthm}
\usepackage{mathrsfs}
\usepackage{tikz-cd}
\usepackage{enumitem}

\frenchsetup{AutoSpacePunctuation=false}
\tikzcdset{ampersand replacement=\&}

\newcommand{\printedpage}[1]{\marginpar{\footnotesize p.~#1}}
\newcommand{\msp}[1]{\marginpar{\footnotesize [#1]}}
\newcommand{\Ho}{\mathrm{Ho}}
\newcommand{\Fib}{\mathrm{Fib}}

\begin{document}

\printedpage{165}
% head: A. Grothendieck : Dérivateurs, Ch. XII, Édition des Universités de Montpellier II et Paris VII
% foot: Version 30 mars 2003

\textit{La loi fonctorielle s'obtient à l'aide du lemme 1 p. 249, qui permet de se borner aux arguments $\begin{tikzcd}[ampersand replacement=\&] X_1 \arrow[d, "p_1"'] \& \\ S \& X_2 \arrow[l, "p_2"] \end{tikzcd}$ tels que $p_1,p_2\in \mathit{Fib}_W$, de sorte que l'on trouve un foncteur}
\[
\begin{tikzcd}[ampersand replacement=\&]
X_1 \arrow[d] \& \\
S \& X_2 \arrow[l]
\end{tikzcd}
\qquad\longmapsto\qquad
\begin{tikzcd}[ampersand replacement=\&]
X_1 \arrow[d] \& X_1\times_S X_2 \arrow[l] \arrow[d] \\
S \& X_2, \arrow[l]
\end{tikzcd}
\]
\textit{compatible aux localiseurs, et passant aux catégories de fractions} $\Ho_W(\Psi)$ \textit{et} $\Ho_W(\tilde{\Psi})$ \textit{respectivement.}

\textbf{NB} En procédant de façon similaire, plus simple (sans besoin d'un recours au lemme 1), on trouve l'existence d'un \textit{adjoint à gauche} $i_!$ de $i^* : \Ho_W(\Psi)\longrightarrow \Ho_W(\tilde{\Psi})$, sous réserve de l'existence d'un objet initial $\varphi$ dans $\mathcal{M}$. Cet foncteur $i_!$ est donné par % [sic] Cet foncteur -> Ce foncteur, "foncteur" begins with a consonant

\msp{page 255}
\[
i_!\left(
\begin{tikzcd}[ampersand replacement=\&]
X_1 \arrow[d] \& \\
S \& X_2 \arrow[l]
\end{tikzcd}
\right)
\qquad = \qquad
\left(
\begin{tikzcd}[ampersand replacement=\&]
X_1 \arrow[d] \& \varphi \arrow[l] \arrow[d] \\
S \& X_2 \arrow[l]
\end{tikzcd}
\right),
\]
la loi fonctorielle étant évidente.

Ce n'est pas encore la fin du sorite! Il faut encore interpréter le foncteur composé
\[
\Ho_W(\Psi)\xrightarrow{\ i_*\ }\Ho_W(\tilde{\Psi})\xrightarrow{\ \varepsilon_3^*\ }\Ho_W(e)=\Ho_W,
\]
où
\[
\varepsilon_3 : e\xrightarrow{\ \sim\ }\{3\}
\qquad\longrightarrow\qquad
\begin{tikzcd}[ampersand replacement=\&, column sep=small, row sep=small]
\& 1 \arrow[dl] \& \\
3 \& \& 0 \arrow[ul] \arrow[dl] \\
\& 2 \arrow[ul] \&
\end{tikzcd}
\]
est l'inclusion du sommet 3 dans $\tilde{\Psi}$. Si $\Ho_W(\Psi)$ est interprété, conformément au lemme 1, comme $\mathcal{F}W_{\mathcal{F}}^{-1}$, ce foncteur composé n'est autre que
\[
\begin{tikzcd}[ampersand replacement=\&]
X_1 \arrow[d] \& \\
S \& X_2 \arrow[l]
\end{tikzcd}
\qquad\longmapsto\qquad X_1\times_S X_2
\]
\printedpage{166}
% head: A. Grothendieck : Dérivateurs, Ch. XII, Édition des Universités de Montpellier II et Paris VII
% foot: Version 30 mars 2003
(pour $\left(\begin{tikzcd}[ampersand replacement=\&] X_1 \arrow[d, "p_1"'] \& \\ S \& X_2 \arrow[l, "p_2"] \end{tikzcd}\right)\in \mathrm{Ob}\,\mathcal{F}$, i.e. $p_1,p_2\in\Fib_W$.)

On veut interpréter ce foncteur comme un $\varprojlim$ $W$-homotopique, i.e. vu comme un `produit fibré $W$-homotopique'. De façon précise, on a ceci :

\msp{page 256}
\textbf{Corollaire 1.} \textit{Considérons $p_\Psi : \Psi\longrightarrow e$, et le foncteur}
\[
p_\Psi^* : \Ho_W(e)=\Ho_W\longrightarrow\Ho_W(\Psi),
\]
\textit{déduit par passage aux catégories de fractions du foncteur}
\[
\begin{array}{ccc}
\mathcal{M} & \longrightarrow & \mathcal{M}(\Psi) \\
X & \longrightarrow & \left(
\begin{tikzcd}[ampersand replacement=\&]
X \arrow[d] \& \\
X \& X \arrow[l]
\end{tikzcd}
\right).
\end{array}
\]
\textit{Ce foncteur admet le foncteur $\varepsilon_3^*i_*$ précédent comme adjoint à droite, i.e.}
\[
p_{\Psi\,*}=\varepsilon_3^*i_* : \Ho_W(\Psi)\longrightarrow\Ho_W.
\]

\textsc{Démonstration}. On peut regarder $p_\psi$ comme le composé % [sic] p_\psi -> p_\Psi, the subscript is the capital \Psi everywhere else
\[
\Psi\xrightarrow{\ i\ }\tilde{\Psi}\xrightarrow{\ p_{\tilde{\Psi}}\ }e,
\]
et on a donc
\[
p_\Psi^*\simeq i^*(p_{\tilde{\Psi}})^*.
\]
Donc l'existence et la détermination d'un adjoint à droite pour $p_\Psi^*$ résultera de l'existence et de la détermination d'adjoints à droite $i_*$, $(p_{\tilde{\Psi}})_*$ de $i^*$, $(p_{\tilde{\Psi}})^*$, via
\[
(p_\Psi)_*\simeq(p_{\tilde{\Psi}})_*i_*.
\]

\msp{page 257}
Donc le corollaire 1 résulte de la formule
\[
p_{\tilde{\Psi}\,*}=\varepsilon_3^*,
\]
qui elle-même est un cas particulier du

\printedpage{167}
% head: A. Grothendieck : Dérivateurs, Ch. XII, Édition des Universités de Montpellier II et Paris VII
% foot: Version 30 mars 2003

\textbf{Lemme.} \textit{Soit $I$ dans} Cat\textit{, avec objet final $e$, considérons l'inclusion}
\[
\varepsilon : e\longrightarrow I,\qquad \textit{et } p : I\longrightarrow e.
\]
\textit{Soit $(\mathcal{M},W)$ une catégorie de modèles, et considérons}
\[
\Ho_W(I)\xleftarrow{\ p^*\ }\Ho_W.
\]
\textit{Ce foncteur admet un adjoint à gauche} \texttt{[plutôt à droite]} \textit{$p_*$, lequel est isomorphe à $\varepsilon^*$.}

\textsc{Démonstration}. Considérons le couple de foncteurs adjoints
\[
\begin{tikzcd}[ampersand replacement=\&, column sep=large]
\mathcal{M}(I) \arrow[r, shift right, "p_*=\varprojlim\limits_I"'] \& \mathcal{M} \arrow[l, shift right, "p^*"']
\end{tikzcd}\ ;
\]
comme $I$ admet l'objet final $e$, le foncteur $p_*$ est isomorphe à $\varepsilon^*$. Ces foncteurs $p^*$, $\varepsilon^*$ sont compatibles aux localiseurs $W(I)$, $W$, donc on conclut par le lemme page 248.

Il ne reste plus qu'à prouver les lemmes 1 et 2. Je le ferai demain.

\msp{page 258}
Aujourd'hui, j'ai envie encore de donner les sorites des carrés $W$-cartésiens dans une

\textbf{Proposition 19 :} \textit{Soit $\begin{tikzcd}[ampersand replacement=\&] X_1 \arrow[d, "p_1"'] \& X_3 \arrow[l, "f_1"'] \arrow[d, "f_2"] \\ S \& X_2 \arrow[l, "p_2"] \end{tikzcd}$ un carré commutatif dans $\mathcal{M}$. Considérons des factorisations de $p_1$, $p_2$ en}
\[
X_1\xrightarrow{\ i_1\ }\bar{X}_1\xrightarrow{\ \bar{p}_1\ }S,\qquad X_2\xrightarrow{\ i_2\ }\bar{X}_2\xrightarrow{\ \bar{p}_2\ }S,
\]
\textit{$i_1,i_2\in W$, $\bar{p}_1,\bar{p}_2\in F_W$} \texttt{[$=$ Fib$_W$]}\textit{, et considérons les flèches induites par $X_3\longrightarrow X_1\times_S X_2$ (déduite de $Q$) via $\begin{tikzcd}[ampersand replacement=\&, row sep=small] \& \bar{X}_1\times_S X_2 \\ X_1\times_S X_2 \arrow[ur] \arrow[r] \arrow[dr] \& \bar{X}_1\times_S\bar{X}_2 \\ \& X_1\times_S\bar{X}_2 \end{tikzcd}$, de sorte qu'on trouve un diagramme commutatif}
\[
\begin{tikzcd}[ampersand replacement=\&, row sep=small]
\& \bar{X}_1\times_S X_2 \arrow[d, "j_2"'] \\
X_3 \arrow[ur, "\pi_1"] \arrow[r, "\pi_0"] \arrow[dr, "\pi_2"] \& \bar{X}_1\times_S\bar{X}_2 \\
\& X_1\times_S\bar{X}_2, \arrow[u, "j_1"]
\end{tikzcd}
\]
\textit{où $j_1=i_1\times_S\mathrm{id}_{\bar{X}_2}$ et $j_2=\mathrm{id}_{\bar{X}_1}\times i_2$ sont dans $W$ (car $i_1$ et $i_2$ le sont, et $\bar{X}_1$ et $\bar{X}_2$ sont dans $\underline{\mathrm{Fib}}_W S$). On a donc}
\[
(\pi_1\in W)\qquad\iff\qquad(\pi_2\in W)\qquad\qquad(\iff\ \pi_0\in W),
\]
\textit{ce qui prouve que la condition $\pi_1\in W$, qui semble dépendre du choix de la factorisation par $i_1$, n'en dépend pas, et idem pour $\pi_2\in W$, et pour $\pi_0\in W$ (qui semble dépendre à la fois de $i_1$, et de $i_2$).}

\printedpage{168}
% head: A. Grothendieck : Dérivateurs, Ch. XII, Édition des Universités de Montpellier II et Paris VII
% foot: Version 30 mars 2003

\msp{page 259}
\textbf{Définition.} \textit{Quand ces conditions sont satisfaites, on dit que le carré $Q$ est $W$-cartésien.}

\textbf{Corollaire 1.} \textit{Si $Q$ est $W$-cartésien, alors l'image de $Q$ dans} $\Ho_W(\tilde{\Psi})$ \textit{est dans l'image essentielle de}
\[
i_* : \Ho_W(\Psi)\longrightarrow\Ho_W(\tilde{\Psi}),
\]
\textit{i.e. la flèche canonique dans} $\Ho_W(\tilde{\Psi})$
\[
Q\longrightarrow i_*i^*Q
\]
\textit{est} \texttt{[un]} \textit{isomorphisme. La réciproque est vraie si $W$ est fortement saturé.}

\textbf{Corollaire 2.} \textit{Considérons un diagramme commutatif dans $\mathcal{M}$}
\[
\begin{tikzcd}[ampersand replacement=\&]
X_0' \arrow[d] \& X_1' \arrow[l] \arrow[d] \arrow[dl, phantom, "\textit{1}"] \& X_2' \arrow[l] \arrow[d] \arrow[dl, phantom, "\textit{2}"] \\
X_0 \& X_1 \arrow[l] \& X_2, \arrow[l]
\end{tikzcd}
\tag*{(D)}
\]
\textit{d'où deux carrés commutatifs 1, 2, et un carré composé}
\[
\begin{tikzcd}[ampersand replacement=\&]
X_0' \arrow[d] \& X_2' \arrow[l] \arrow[d] \arrow[dl, phantom, "\textit{3}"] \\
X_0 \& X_2. \arrow[l]
\end{tikzcd}
\]
\textit{Supposons le carré 1 $W$-cartésien. Alors le carré 2 est $W$-cartésien si et seulement si le carré 3 l'est.}

\textsc{Démonstration}. Considérons une factorisation de $\begin{tikzcd}[ampersand replacement=\&, row sep=small] X_0' \arrow[d] \\ X_0 \end{tikzcd}$ en
\[
\begin{tikzcd}[ampersand replacement=\&, column sep=small, row sep=small]
\bar{X}_0' \arrow[dr] \& X_0' \arrow[l] \arrow[d] \\
\& X_0
\end{tikzcd}
\]
avec $(X_0'\longrightarrow\bar{X}_0')\in W$, et $(\bar{X}_0'\longrightarrow X_0)\in\Fib_W$. On peut alors remplacer dans (D) $X_0'$ par $\bar{X}_0'$, obtenant un diagramme $(\bar{D})$, et on voit aussitôt
\printedpage{169}
% head: A. Grothendieck : Dérivateurs, Ch. XII, Édition des Universités de Montpellier II et Paris VII
% foot: Version 30 mars 2003
\msp{page 260}
que chacun des carrés 1(D), 2(D), 3(D) est $W$-cartésien si et seulement si le carré correspondant $1(\bar{D})$, $2(\bar{D})$ ou $3(\bar{D})$ l'est. Cela montre qu'on peut supposer $X_0\longrightarrow S$ dans $\Fib_W S$. Considérons alors le diagramme commutatif déduit de (D)
\[
\begin{tikzcd}[ampersand replacement=\&, column sep=small]
X_0' \arrow[d] \& \tilde{X}_1'\overset{\text{déf}}{=}X_0'\times_{X_0}X_1 \arrow[l] \arrow[d] \& X_1' \arrow[l, "\alpha"'] \& X_2' \arrow[l] \arrow[d] \\
X_0 \& X_1 \arrow[l] \& \& X_2, \arrow[ll]
\end{tikzcd}
\]
où $\alpha\in W$ en vertu de l'hypothèse $W$-cartésienne sur le carré 1. On voit encore que le carré 2 est $W$-cartésien si et seulement si le carré correspondant avec $X_1'$ remplacé par $\tilde{X}_1'$ l'est, quant aux carrés 1, 3, ils ne changent pas par cette modification. Cela montre qu'on peut supposer le carré 1 \textit{cartésien.} Considérons alors le diagramme
\[
\begin{tikzcd}[ampersand replacement=\&]
X_0' \arrow[d, "\mathrm{Fib}_W"'] \& X_1' \arrow[l] \arrow[d, "\mathrm{Fib}_W"] \arrow[dl, phantom, "\text{cart.}"] \& \tilde{X}_2'\overset{\text{déf}}{=}X_1'\times_{X_1}X_2 \arrow[l] \arrow[d, "\mathrm{Fib}_W"] \arrow[dl, phantom, "\text{cart.}"] \& X_2' \arrow[l, "\beta"'] \\
X_0 \& X_1 \arrow[l] \& X_2. \arrow[l]
\end{tikzcd}
\]
Alors la proposition 19 et la définition montrent que le carré 2 est $W$-cartésien si et seulenment si $\beta\in W$, et de même pour le carré 3. Ce qui achève la démonstration. % [sic] seulenment -> seulement, misspelling

\textbf{Corollaire 2.} \textit{Soit $\begin{tikzcd}[ampersand replacement=\&] X_1 \arrow[d, "p_1"'] \& X_3 \arrow[l, "f_2"'] \arrow[d, "f_1"] \\ S \& X_2 \arrow[l, "p_2"] \end{tikzcd}$ un carré $W$-cartésien. Alors}
\[
\begin{aligned}
p_1\in W\quad&\Longrightarrow\quad f_1\in W\\
p_2\in W\quad&\Longrightarrow\quad f_2\in W,
\end{aligned}
\]
\textit{plus précisément, si $p_i\in W$, alors $Q$ est $W$-cartésien si et seulement si $f_i\in W$.}

\msp{page 261}
Supposons $p_1\in W$. Alors dans la factorisation de $p_1$ en $X_1\longrightarrow\bar{X}_1\longrightarrow S$, on peut prendre $\bar{X}_1=S$. Donc le carré est $W$-cartésien si et seulement si
\[
(X_3\longrightarrow\bar{X}_1\times_S X_2=X_2)\in W,
\]
i.e. $f_1\in W$, q.e.d.

Bien sûr, on ne peut en général affirmer les implications inverses : $f_i\in W$ n'implique nullement $p_i\in W$.

\end{document}
